%============================================================
% Callout Box: Response Immersion Theorem
%============================================================
\begin{tcolorbox}[
colback=blue!3,
colframe=blue!60!black,
title=\textbf{Response Immersion Theorem: Why It Matters},
fonttitle=\bfseries]
A central result of this work is that the reciprocal response metric
$G_{\mu\nu}$ and the antisymmetric response curvature
$F_{\mu\nu}$ need not be regarded as independent geometric objects.
Whenever both arise from the same underlying response immersion,
they satisfy the exact identity
\begin{equation}
\boxed{
\det G = F_{12}^{\,2}.
}
\end{equation}
This relation has an important physical consequence.  The metric,
which characterizes reciprocal susceptibility and
distinguishability, and the curvature, which characterizes cyclic
transport and geometric work, become complementary projections of
the same underlying response geometry rather than unrelated
quantities.
Equivalently, the response operator
\[
J = G^{-1}F
\]
defines a compatible complex structure satisfying
\[
J^2=-I,
\]
so that the local response manifold possesses a K\"ahler geometry.
Thus, the appearance of K\"ahler structure is not imposed as an
additional mathematical assumption, but emerges naturally whenever
the reciprocal and nonreciprocal sectors of response share a common
geometric origin.
The theorem therefore provides a unifying geometric interpretation
of response theory: \emph{reciprocal response determines local
distance, nonreciprocal response determines local circulation, and
together they define the intrinsic geometry of the response
manifold.}
\end{tcolorbox}